\chapter{Postpartum Depression Screening}
\section*{What Is This Test?} Postpartum depression screening identifies depressive and anxiety symptoms during pregnancy or after birth.\section*{Alternative Names} Edinburgh Postnatal Depression Scale; EPDS; postpartum mood screening.\section*{Principle} Validated questionnaires are combined with clinical interview and assessment of safety, functioning, sleep, and support.\section*{Why This Test Is Ordered} It is used when mood, interest, bonding, sleep, worry, or functioning changes during the perinatal period.\section*{Primary vs Secondary Test} Screening is preliminary; diagnosis and treatment planning require clinical evaluation.\section*{How This Test Is Performed} The patient completes questions and discusses symptoms, intrusive thoughts, mania, psychosis, substances, and support.\section*{What Preparation Is Needed} None. Report thoughts of self-harm or harm to the baby immediately.\section*{Understanding Results} A positive screen requires prompt assessment; postpartum blues are common but persistent, severe, or dangerous symptoms need care.\section*{Follow-up Tests} Psychiatric assessment, thyroid or medical testing when indicated, and urgent safety evaluation may follow.\section*{References} \begin{enumerate}\item MedlinePlus. Postpartum Depression Screening.\item American College of Obstetricians and Gynecologists. Perinatal mental-health screening guidance.\end{enumerate}\clearpage\section*{Understanding Results and Follow-up}
The laboratory report should identify the specimen, method, units, reference interval or decision limit, and whether the result is positive, negative, abnormal, or inconclusive. Reference intervals vary with age, sex, pregnancy, specimen type, assay, and laboratory; a result outside a range is not automatically a diagnosis, and a result within a range does not always exclude disease. Discuss unexpected results with the ordering clinician, who can decide whether repeat or confirmatory testing, treatment, or referral is appropriate. Do not change medicines or delay urgent care based on an isolated result.
\section*{Regulatory Status and Authorization Date}This chapter describes a test category, not a specific commercial product. FDA, CDSCO, EU IVDR, and MHRA approval or registration dates therefore cannot be assigned without the assay manufacturer, product name, instrument, and intended use. For laboratory-developed tests, an individual agency approval date may not exist. See the Preface for the interpretation of regulatory status.
\clearpage
